\section{Parts}

\begin{tabular}{| l | c | l | l |}
\hline
Part & Icon & Base Price & Base Tech \\
\hline
Control Node & \includegraphics[scale=0.70]{images/part-controlnode.png} & & \\
Flux Modulator & \includegraphics[scale=0.70]{images/part-fluxmodulator.png} & &\\
Photon Bridge & \includegraphics[scale=0.70]{images/part-photonbridge.png} & &\\
ST Distorter & \includegraphics[scale=0.70]{images/part-stdistorter.png} & &\\
Cryofuse & \includegraphics[scale=0.70]{images/part-cryofuse.png} & &\\
Plasma Tube & \includegraphics[scale=0.70]{images/part-plasmatube.png} & &\\
Scanner & \includegraphics[scale=0.70]{images/part-scanner.png} & &\\
J Junc Module & \includegraphics[scale=0.70]{images/part-jjunc.png} & &\\
\hline
\end{tabular}

\subsubsection{Control Node}
The SunDog needs control nodes to function correctly.  Control nodes are
used in each system row in each system bay and can not be shunted, so it's
necessary to have a reasonable amount of them on hand to keep the ship
running smoothly.  

\subsubsection{Flux Modulator}
Flux Modulators are used with the shields, sub-c engines and warp drives.

\subsubsection{Photon Bridge}
Photon bridges are used with tactical and the warp engines.

\subsubsection{ST Distorter}
Space/time distorters (ST Distorters) are used only with the warp and
sub-light (sub-c) engines.

\subsubsection{Cryofuse}

Open Issues:
\begin{itemize}
\item What exactly is a cryofuse?  Cryogenics fuse?  Some kind of cold part?
\end{itemize}

\subsubsection{Plasma Tube}

Plasma tubes are used as part of the cannon and laser systems aboard the
SunDog.

\subsubsection{Scanner}
\subsubsection{J Junc Module}
\subsubsection{Shunt}


\subsection{Special Parts}
A number special parts can be added to the Sun Dog to give it extra
capabilities.

\begin{tabular}{| l | c | l | l |}
\hline
Part & Icon & Base Price & Base Tech \\
\hline
Ground Scanner & \includegraphics[scale=0.70]{images/part-groundscanner.png} & & \\
Cloaker & \includegraphics[scale=0.70]{images/part-cloaker.png} & & \\
Decloaker & \includegraphics[scale=0.70]{images/part-decloaker.png} & & \\
Concentrator & \includegraphics[scale=0.70]{images/part-concentrator.png} & & \\
Booster & & 200,000 & Only found in Banville \\
\hline
\end{tabular}

Open Issues:
\begin{itemize}
\item Do we want to have an auto-slew like the Apple II version?
\item Is there an image for the booster?
\end{itemize}

\subsubsection{Ground Scanner}

The ground scanner gives the SunDog the capability to land at any city
on any planet, regardless of whether there is a space port located in
the city or not.  This capability works both from orbit, and for
city-to-city flights so that the ship does not have to leave the planet's
atmosphere (useful when wanting to avoid pirates).

\subsubsection{Cloaker}

The cloaker can be used to disguise the SunDog from its enemies.  This allows
the ship to evade cannon and laser fire from attacking ships until the
SunDog can land or warp to a different solar system.  Using the cloaker
drains a tremendous amount of fuel and will automatically shut off if there
is not enough fuel remaining to power it.

Open Issues:
\begin{itemize}
\item What is the rate of fuel drain for using the cloaker?
\end{itemize}

\subsubsection{Decloaker}

\subsubsection{Concentrator}

The concentrator allows the cannon and laser weapons to inflict more damage
than the standard system.

Open Issues:
\begin{itemize}
\item Is there a fuel penalty for using the concentrator?
\end{itemize}

\subsubsection{Booster}

The booster allows the SunDog to warp further distances than it normally
can travel.  This is the only way the ship can make it to Enlie in order
to pick up the remaining cryogens to complete Phase 9 of Banville (see the
Banville section).

Open Issues:
\begin{itemize}
\item When does the booster become available?
\end{itemize}


